% Generado por scripts/figuras/tablas.py - no editar a mano
\begin{table}[htbp]
  \centering
  \caption{Recall del ataque antes y despues de camuflar los primeros bytes con los de un cliente legitimo.}
  \label{tab:evasión}
  \small
  \begin{tabular}{llrrr}
    \toprule
    \textbf{Escenario} & \textbf{Vista} & \textbf{Original} & \textbf{Evadido} & \textbf{Caida} \\
    \midrule
    SSH cifrado (14-02) & payload (CNN) & 1.0000 & 0.3538 & \textbf{0.6462} \\
    SSH cifrado (14-02) & conducta (ráfaga) & 1.0000 & 1.0000 & invariante \\
    DoS volumetrico (16-02, 200 B) & payload (CNN) & 0.7242 & 0.4258 & \textbf{0.2984} \\
    DoS volumetrico (16-02, 200 B) & conducta (ráfaga) & 0.9992 & 0.9992 & invariante \\
    \bottomrule
  \end{tabular}
  \par\smallskip\footnotesize{Se alteran unicamente los primeros 48 bytes (SSH) o 200 bytes (DoS) de cada flujo de ataque; el resto del flujo no cambia. La vista conductual no mira el contenido, por lo que su recall no varia.}
\end{table}
